\chapter{Marco teórico y estado del arte}

Este capítulo reúne los conceptos necesarios para seguir el resto del documento.
Se parte de los fundamentos de la síntesis de sonido, se detalla la técnica
concreta sobre la que se construye el instrumento y se introducen los
autoencoders y la noción de espacio latente. El capítulo termina con una
revisión de los trabajos previos que aplican redes neuronales a la generación de
formas de onda, que es donde se sitúa la aportación de este trabajo.

\section{Fundamentos de síntesis de sonido}

Un sonido periódico queda descrito por tres magnitudes: la frecuencia
fundamental, que determina la altura de la nota; la amplitud, que determina su
intensidad; y el contenido armónico, que determina el timbre. Dos instrumentos
que tocan la misma nota con la misma intensidad se distinguen precisamente por
la distribución de energía entre los armónicos de esa fundamental. Estos armónicos
dependen de características físicas del instrumento, como la forma de la caja de resonancia
o la rigidez de las cuerdas. En la sintetización de sonido, en cambio, el timbre 
se genera mediante algoritmos que producen la señal de audio directamente, sin 
necesidad de un instrumento físico. 

La síntesis de sonido agrupa las técnicas que permiten generar señales de audio
por medios electrónicos. Las más habituales son las siguientes
\cite{Wilson2013}:

\begin{itemize}
    \item \textbf{Síntesis aditiva.} Construye el sonido sumando osciladores
          senoidales de distinta frecuencia y amplitud. Ofrece un control muy
          preciso del espectro, a costa de un coste computacional elevado
          cuando el número de frecuencias a sumar es alto.
    \item \textbf{Síntesis sustractiva.} Parte de una señal rica en armónicos y
          emplea filtros para eliminar las frecuencias no deseadas. Es la base
          de los sintetizadores analógicos clásicos y sigue vigente en los
          digitales.
    \item \textbf{Síntesis por modulación de frecuencia.} Modula la frecuencia
          de un oscilador con la salida de otro, generando espectros complejos
          con muy pocos operadores.
    \item \textbf{Síntesis por tabla de ondas.} Almacena en memoria uno o varios
          ciclos de forma de onda y los reproduce leyendo la tabla a la
          velocidad adecuada. Es la técnica sobre la que se construye este
          trabajo y se detalla en el apartado siguiente.
\end{itemize}

Con independencia de la técnica empleada, un sonido musical necesita evolucionar
en el tiempo. Una nota que se mantuviese con amplitud constante desde que se
pulsa la tecla hasta que se suelta resultaría artificial. Esa evolución se
controla mediante una envolvente, y la más extendida es la \gls{adsr}, que
divide la vida de la nota en cuatro tramos:
\begin{itemize}
      \item \textbf{Ataque (Attack).} El tiempo que tarda la señal en alcanzar su nivel máximo.
      \item \textbf{Decaimiento (Decay).} El tiempo que tarda en descender el nivel desde el pico máximo hasta el nivel de sostenimiento.
      \item \textbf{Sostenimiento (Sustain).} El nivel que se mantiene mientras la tecla siga pulsada.
      \item \textbf{Liberación (Release).} El tiempo que tarda en devolver la señal a cero cuando se suelta la tecla.
  \end{itemize}

La Figura~\ref{fig:adsr} muestra la forma que adopta la envolvente y cómo se
reparten esos cuatro tramos a lo largo de la duración de la nota.

\begin{figure}[htbp]
    \centering
    \includegraphics[width=0.85\textwidth]{adsr.pdf}
    \caption{Envolvente ADSR. Los tres primeros tramos dependen de cuánto tiempo
             se mantenga pulsada la tecla; la liberación empieza al soltarla.}
    \label{fig:adsr}
\end{figure}

El otro elemento habitual en la cadena de un sintetizador es el filtro. Un
\gls{lpf} atenúa las frecuencias por encima de una frecuencia de corte, un
\gls{hpf} hace lo contrario, atenúa las frecuencias por debajo de una frecuencia
de corte; y un \gls{bpf} deja pasar únicamente una banda de frecuencia
intermedia. El parámetro de resonancia (Q) realza la respuesta en torno a la
frecuencia de corte, lo que da al filtro buena parte de su carácter sonoro. 

La Figura~\ref{fig:filtros} compara la respuesta en magnitud de los tres tipos
sobre una misma frecuencia de corte.

\begin{figure}[htbp]
    \centering
    \includegraphics[width=0.85\textwidth]{filtros.pdf}
    \caption{Respuesta en magnitud de un filtro paso bajo, uno paso alto y uno
             paso banda de segundo orden, todos con la frecuencia de corte en
             1~kHz y factor de calidad $Q = 0{,}707$.}
    \label{fig:filtros}
\end{figure}

El orden clásico de la cadena en un sintetizador sustractivo es oscilador, filtro y
amplificador controlado por la envolvente (\gls{adsr}), y es el que se ha mantenido en este
trabajo.

\section{Síntesis por tabla de ondas}

La síntesis wavetable se apoya en una observación sencilla. Si una señal es
periódica, basta con almacenar un único ciclo y repetirlo. En lugar de evaluar
una función trigonométrica para cada muestra, se guarda ese ciclo en una tabla
de $N$ posiciones y se recorre cíclicamente a la velocidad que corresponda a la
nota deseada \cite{BristowJohnson1996}.

El mecanismo que gobierna esa lectura es el acumulador de fase. Se trata de una
variable que avanza una cantidad fija en cada muestra y que vuelve al principio
al llegar al final de la tabla. Siendo $f_0$ la frecuencia de la nota, $f_s$ la
frecuencia de muestreo y $N$ el tamaño de la tabla, el incremento de fase por
muestra es el de la Ecuación~\eqref{eq:incfase}.

\begin{equation}
\label{eq:incfase}
\Delta\varphi = \frac{f_0 \cdot N}{f_s}
\end{equation}

El acumulador avanza según la Ecuación~\eqref{eq:acumulador}, donde el módulo
implementa el retorno al inicio de la tabla al completarse cada ciclo.

\begin{equation}
\label{eq:acumulador}
\varphi[n+1] = \left( \varphi[n] + \Delta\varphi \right) \bmod N
\end{equation}

Como $\Delta\varphi$ no es en general un número entero, el acumulador apunta
casi siempre a una posición intermedia entre dos muestras de la tabla. La
solución habitual es interpolar linealmente entre la muestra anterior y la
siguiente, tal como recoge la Ecuación~\eqref{eq:interplineal}, donde
$i = \lfloor \varphi \rfloor$ y $\alpha = \varphi - i$ es la parte fraccionaria.

\begin{equation}
\label{eq:interplineal}
s[n] = (1-\alpha) \cdot T[i] + \alpha \cdot T[(i+1) \bmod N]
\end{equation}

La Figura~\ref{fig:wavetable} ilustra el proceso completo. El acumulador señala
una posición fraccionaria dentro de la tabla y el valor de salida se obtiene
pesando las dos muestras que lo rodean.

\begin{figure}[htbp]
    \centering
    \includegraphics[width=\textwidth]{wavetable.pdf}
    \caption{Lectura de una tabla de ondas. A la izquierda, la tabla almacenada y
             la posición actual del acumulador de fase $\varphi$. A la derecha,
             el detalle de la interpolación lineal entre las dos muestras
             vecinas.}
    \label{fig:wavetable}
\end{figure}

El coste de generar una muestra se reduce así a una suma, dos multiplicaciones y
un acceso a memoria, lo que hace la técnica especialmente adecuada para
microcontroladores. Bristow-Johnson señala además que, para tonos musicales
cuasiperiódicos, este método resulta equivalente a la síntesis aditiva pero
exigiendo mucho menos cálculo en tiempo real \cite{BristowJohnson1996}, que es
justo el compromiso que interesa en un sistema embebido. A cambio, el timbre
queda fijado por el contenido de la tabla, de modo que cambiar de sonido implica
cambiar de tabla.

\subsection{Aliasing y limitación de banda}

Una tabla de $N$ muestras puede contener hasta $N/2$ armónicos. Al reproducirla
a frecuencias altas, los armónicos superiores acaban rebasando la frecuencia de
Nyquist y se reflejan hacia el interior del espectro como componentes de
frecuencia no armónica, audibles como un timbre metálico y desafinado. Este
fenómeno se conoce como \emph{aliasing} y es el principal problema de la
reproducción de tablas a frecuencias elevadas.

La solución habitual consiste en almacenar varias versiones de cada tabla,
filtradas para distintos rangos de nota, y seleccionar en cada momento la que
corresponda al tono que se está tocando \cite{Valimaki2007}. Conviene señalar
que este problema es independiente del tamaño de la tabla y del proceso de
remuestreo empleado para construirla. Aumentar $N$ no elimina el
\emph{aliasing}, solo cambia cuántos armónicos hay disponibles para reflejarse.
Por las razones que se exponen en el capítulo~\ref{cap:diseno}, la limitación de
banda queda fuera del alcance de este trabajo.

\section{Autoencoders y espacio latente}

Un autoencoder es una red neuronal entrenada para reproducir su propia entrada.
Está formado por dos bloques. El codificador transforma la entrada en una
representación de dimensión mucho menor, y el decodificador intenta reconstruir
la entrada original a partir de esa representación. La representación intermedia
recibe el nombre de código o vector latente, y el espacio en el que vive, espacio
latente.

El interés de esta arquitectura no está en la reconstrucción en sí, que sería
trivial si no se impusiera ninguna restricción, sino en el cuello de botella. Al
obligar a que toda la información pase por un número reducido de valores, la red
se ve forzada a descartar lo accesorio y quedarse con los rasgos que de verdad
distinguen unas entradas de otras. Aplicado a un conjunto de formas de onda, el
codificador acaba situando cerca las que comparten características espectrales,
sin que se le haya indicado en ningún momento qué es un armónico.

\subsection{El autoencoder variacional}

Un autoencoder convencional organiza el espacio latente sin ninguna garantía
sobre su estructura. Nada impide que queden zonas vacías entre los puntos
ocupados, y un punto tomado de una de esas zonas puede decodificarse en algo sin
sentido. Para el uso que se pretende aquí, en el que el usuario recorre el plano
con el dedo y pasa continuamente por posiciones que no corresponden a ninguna
onda del conjunto de entrenamiento, esa falta de garantías es un problema.

El autoencoder variacional o \gls{vae} \cite{Kingma2013} resuelve esa limitación
cambiando la naturaleza del código. En lugar de producir un vector de valores
concretos, el codificador produce los parámetros de una distribución de
probabilidad, es decir, un vector de medias $\boldsymbol{\mu}$ y otro de
varianzas $\boldsymbol{\sigma}^2$. El código que recibe el decodificador se
obtiene muestreando esa distribución. Ese muestreo plantea un problema, porque
una operación aleatoria no es derivable y el gradiente no puede atravesarla
durante el entrenamiento. La solución que proponen Kingma y Welling es
reescribir la variable latente como una transformación determinista de la media
y la desviación típica sobre una muestra de ruido tomada aparte, tal como recoge
la Ecuación~\eqref{eq:reparam}. Con esa reformulación, conocida como truco de
reparametrización, la aleatoriedad queda confinada en un término que no depende
de los parámetros de la red, y el gradiente sí puede propagarse.

\begin{equation}
\label{eq:reparam}
\mathbf{z} = \boldsymbol{\mu} + \boldsymbol{\sigma} \odot \boldsymbol{\epsilon},
\qquad \boldsymbol{\epsilon} \sim \mathcal{N}(\mathbf{0}, \mathbf{I})
\end{equation}

A la función de pérdida se le añade además un segundo término que penaliza la
divergencia de Kullback-Leibler entre la distribución producida por el
codificador y una normal estándar. Ese término tira de todas las
distribuciones hacia el mismo punto del espacio, con lo que las zonas ocupadas
se solapan y no quedan huecos. La pérdida total combina entonces el error de
reconstrucción y esa divergencia, como indica la Ecuación~\eqref{eq:vaeloss}.

\begin{equation}
\label{eq:vaeloss}
\mathcal{L} = \underbrace{\mathrm{MSE}(\mathbf{x}, \hat{\mathbf{x}})}_{\text{reconstrucción}}
+ \; \beta \cdot \underbrace{D_{KL}\!\left( q(\mathbf{z}|\mathbf{x}) \,\|\, \mathcal{N}(\mathbf{0},\mathbf{I}) \right)}_{\text{regularización}}
\end{equation}

Para un latente de $d$ dimensiones, la divergencia admite la forma cerrada de la
Ecuación~\eqref{eq:kl}, que es la que se implementa en la práctica.

\begin{equation}
\label{eq:kl}
D_{KL} = -\frac{1}{2} \sum_{j=1}^{d} \left( 1 + \log \sigma_j^2 - \mu_j^2 - \sigma_j^2 \right)
\end{equation}

El factor $\beta$ que aparece en la Ecuación~\eqref{eq:vaeloss} pondera la
importancia relativa de los dos términos y da nombre a la variante conocida como
$\beta$-VAE \cite{Higgins2017}. Higgins y sus coautores lo introducen como un
hiperparámetro que regula el equilibrio entre la capacidad del canal latente y
la precisión de la reconstrucción, y observan que valores mayores que la unidad
empujan al modelo hacia representaciones más ordenadas a costa de reconstruir
peor. Su ajuste supone entonces un compromiso. Valores altos producen un espacio
latente muy regular pero reconstrucciones pobres, porque el modelo prioriza
acercar todas las distribuciones a la normal estándar sobre reproducir fielmente
cada entrada. Valores bajos mejoran la reconstrucción a costa de un espacio
menos continuo.

Existe además un caso degenerado, conocido como colapso posterior, en el que el
modelo opta por ignorar el latente por completo y el término de divergencia se
anula. Bowman y sus coautores se encontraron con este comportamiento al aplicar
un autoencoder variacional a texto, y lo resolvieron introduciendo el peso de la
divergencia de forma progresiva. Durante las primeras épocas el término se anula
o se mantiene muy bajo, de modo que el modelo se centra en aprender a
reconstruir, y solo después se sube hasta su valor definitivo
\cite{Bowman2016}. Es la misma técnica que se ha empleado en este trabajo, tal
como se detalla en el capítulo~\ref{cap:ml}.

\section{Estado del arte}

La aplicación de redes neuronales a la síntesis de audio ha producido en los
últimos años varias líneas de trabajo. Interesa aquí distinguir dos enfoques
distintos según qué es lo que la red genera.

\subsection{Modelos que generan la señal de audio}

Un primer grupo de trabajos emplea la red para generar directamente la forma de
onda de salida, muestra a muestra. El sistema DDSP \cite{Engel2020} combina
elementos clásicos de procesado de señal, como osciladores armónicos y filtros,
con redes neuronales que controlan sus parámetros, de modo que el conocimiento
sobre cómo se genera el sonido queda incorporado a la arquitectura en lugar de
tener que aprenderse desde cero. RAVE \cite{Caillon2021} sigue una vía
diferente. Se trata de un autoencoder variacional entrenado sobre audio en
bruto, con un procedimiento de entrenamiento en dos fases, primero de
aprendizaje de la representación y después de ajuste adversario. Caillon y
Esling añaden un análisis posterior del espacio latente que permite decidir el
compromiso entre fidelidad de reconstrucción y compacidad de la representación,
un problema que aparece de nuevo en este trabajo al fijar la dimensión del
latente. Su interés para este trabajo es que demuestra
que la síntesis neuronal puede alcanzar velocidades de tiempo real, aunque sobre
hardware sensiblemente más potente que un microcontrolador.

\subsection{Modelos que generan tablas de onda}

El segundo grupo, más cercano a lo que se plantea aquí, no genera audio sino
tablas. La red produce una forma de onda de un ciclo que después se reproduce
mediante un oscilador convencional, con lo que el coste de la inferencia queda
desacoplado de la frecuencia de muestreo.

El trabajo de Hantrakul y Yang \cite{Hantrakul2018} es el precedente más
directo. Reutilizan el modelo NSynth WaveNet ya entrenado como autoencoder de
tablas, generando formas de onda de 512 muestras e interpolando entre ellas en
el espacio latente. El resultado se distribuye como \emph{plugin} VST/AU para
estación de trabajo de audio digital.

Dos observaciones suyas resultan especialmente pertinentes aquí. La primera es
de coste. Los propios autores miden unos 4 segundos para codificar y unos 7 para
decodificar una única forma de onda sobre una máquina con GPU, lo que descarta
por completo ejecutar la red durante la reproducción. Su solución fue generar
por adelantado 300 tablas interpolando entre combinaciones de seno, sierra y
triangular, y reproducirlas después con un oscilador convencional. Es exactamente
la misma estrategia que se adopta en este trabajo, y conviene subrayar que allí
la limitación aparece ya sobre una GPU de sobremesa, no sobre un
microcontrolador.

La segunda observación tiene que ver con la estructura del espacio latente.
Hantrakul y Yang comprueban que interpolar a lo largo del segmento que une dos
vectores del espacio produce mezclas utilizables de las dos ondas de partida,
pero que apartarse de ese segmento genera ruido. Dicho de otro modo, en su
modelo no todo el espacio latente es válido, solo lo son las trayectorias entre
puntos conocidos. Para un instrumento en el que el usuario recorre libremente un
plano con el dedo esa garantía no basta, y es una de las razones que justifican
emplear aquí un autoencoder variacional en lugar de uno convencional. Los
propios autores apuntan en esa dirección al plantear como trabajo futuro un
modelo ligero basado precisamente en un autoencoder variacional.

Kreković \cite{Krekovic2021} plantea un oscilador basado en un autoencoder y
estudia el efecto del conjunto de entrenamiento y de los hiperparámetros sobre
la variedad de las formas de onda generadas. Su aportación más pertinente aquí
es la constatación de que emplear las variables latentes directamente como
parámetros de entrada produce una transición suave entre las ondas generadas
cuando los cambios de esos parámetros son pequeños. Es exactamente la propiedad
sobre la que descansa el control táctil de este instrumento.

Más recientemente, Wavespace \cite{Lee2024} propone un generador de tablas
basado en un autoencoder variacional cuyo espacio latente se factoriza por
completo en subespacios asociados a distintos estilos de forma de onda. El
generador se condiciona además con descriptores tímbricos y morfológicos, de
modo que el usuario puede manipular cada subespacio de forma independiente en
lugar de moverse por un plano indiferenciado. Los autores lo prototipan también
como \emph{plugin} de oscilador para estación de trabajo de audio digital, y
señalan que la restricción de velocidad de cálculo sigue siendo el principal
obstáculo cuando se pretende decodificar una forma de onda para cada bloque de
audio.

\subsection{Posicionamiento de este trabajo}

De la revisión anterior se desprende que la generación de tablas de onda
mediante autoencoders y su interpolación en el espacio latente es un problema ya
abordado, y resuelto con solvencia, en varios trabajos previos. La aportación de
este trabajo no está, por tanto, en el modelo.

Lo que ninguno de esos trabajos aborda es la materialización del concepto como
instrumento físico independiente. Todos ellos se presentan como
\emph{plugins}, aplicaciones de escritorio o prototipos de investigación que se
ejecutan sobre un ordenador. La aportación diferencial que se persigue aquí es
la integración embebida completa. Eso supone repartir el sistema entre
microcontroladores de bajo coste, resolver la comunicación entre ellos, mantener
la estabilidad del audio en tiempo real y construir un dispositivo autónomo con
control táctil y entrada \gls{midi} que no necesite un ordenador para funcionar.
El reto no es de aprendizaje automático, sino de ingeniería de sistemas
embebidos.
